\section{Discussion}
\label{sec:discussion}
\subsection{Interpretation of the physical gain}

The main result is that physical-reward training improves measured utility
per draw on the sealed accelerator tasks. Both physical rewards improve over
SFT, with the larger gain obtained by the RF package. This improvement comes
from faster implementations among the outputs that pass the oracle. In
contrast, correctness-only training increases acceptance while reducing
physical utility. The additional-seed extension repeats both directions under
its matched setup flow, and the separate Qwen-Coder experiment shows that the
training recipe can also help another backbone on the earlier task set.

Changing candidate probabilities has practical value because it improves the
expected hardware outcome at a fixed sampling budget. SFT already produces
fast structures for many designs, and optimization makes a smaller set of
passing outputs more likely. The structural summaries and implemented FIR
pair support this interpretation, although they do not show that every gain
arises from the same topology change. The finite samples also leave circuit
novelty unresolved. Likewise, the perfect-selector calculation describes a
separate inference workload and excludes policy-training cost, so it cannot
establish whether reranking or policy training is the better overall strategy.

\subsection{Implications for CAD workflows}

A learned reward is useful when it can distinguish the physical quality of
current policy outputs well enough to guide updates. The RF reward's observed
ranking resolution and sealed efficacy support that role in this domain,
while the archived reward failure and predictor rescoring show why physical
measurements remain necessary. The measured frequency gains also carry costs:
RF uses more LUTs and registers, fewer DSP blocks, and slightly more vectorless
power on common passing support. Pipeline latency can vary within the oracle's
alignment window, so a higher clock frequency does not necessarily reduce
end-to-end latency. A deployment workflow would need explicit resource,
latency, and power constraints, together with a check on general RTL competence
given the matched VerilogEval regression.

The sampling budget and the total compute budget answer different questions.
Matching draws measures the expected outcome of a fixed number of deployment
attempts, but training arms can encounter different numbers of flat groups,
each of which still requires generation and oracle execution. The
perfect-selector comparison also excludes SFT, reward-label, and policy-training
costs. Consequently, the present results do not identify the runtime, energy,
or deployment volume at which policy training becomes more economical than
repeated sampling and physical selection.

\subsection{Scope and validity}

The split evaluates unseen parameters within known streaming-accelerator
generators. Its findings therefore do not establish competence on
control-dominated systems, memories, multiple clocks, or free-form industrial
RTL. Median filters appear only in extrapolation, which prevents us from
separating their family effect from the regime effect. The Qwen study provides
evidence on a second backbone but uses a narrower, earlier task universe.
These scope limits also coexist with unknown overlap between the tested tasks
and base-model pretraining data.

All RF and correctness-only training replications share one SFT initialization.
Design-bootstrap intervals condition on the evaluated checkpoints and draws;
they do not capture full-pipeline training variance. Family cells are small and
unequal. RF versus SFT is the sole prespecified primary contrast; secondary
intervals have no multiplicity adjustment. RF versus MLP changes predictor class,
canonicalization, and features jointly. It supports a reward-package comparison,
while individual component effects and matched external-method superiority
remain untested. In particular, no matched DPO or recent external RTL-optimization
baseline was executed.

The oracle supplies sampled, latency-aligned acceptance at a near-unity
threshold. It does not prove exact functional equivalence. The original physical
endpoint uses one device and a single requested-period flow, with constraints
introduced after synthesis. Its slack variable has a zero default when a valid
setup path is absent; the retained primary rows have not been reclassified under
a stronger contract. The complete coverage diagnostic records one missing-path
case and inconclusive implementations alongside close timing agreement on common
valid candidates. The period-bracketed pilot and additional-seed flow add
stronger setup checks but do not enforce hold or pulse-width closure. The
reverse-FIR pilot outlier also limits absolute-frequency calibration. Broader
device, tool, and operating-corner transfer remains untested.

The RF repair hypothesis was motivated by earlier MLP outcomes. Its sealed
efficacy split stayed unseen until checkpoints and the reward-eligible validation
contract were frozen, but that contract followed an unsuccessful broader gate.
The board subset likewise required feasibility and sweep amendments before
measurement. These histories are retained in the Supplementary Material:
the board result is descriptive, and the repair was prospectively tested after
being motivated by earlier evidence. Contextual API and selected-candidate HLS
comparisons have different reporting objects from the primary policy endpoint.

\subsection{Artifact and claim traceability}

The research archive retains the manifests, sampled RTL and multiplicities, oracle
summaries, implementation failures, board traces, and numerical claim
ledgers needed to trace the reported results.\footnote{
Project repository: \url{https://github.com/MoAmineHallam/FPGA}.}
The accompanying supplementary package provides the worked RTL example and
recomputes the completed RF versus SFT comparison from its archived draws and
physical records, with a file manifest linking the inputs to the analysis.
Regeneration checks design coverage, draw denominators, implementation records,
source hashes, and every numerical claim in the paper and supplement. Each
completed supporting study retains its original evaluation contract. Prepared
protocols, training-only seed logs, the asymmetric historical board comparison,
and the incomplete smaller-student evaluation remain in the archive but do
not enter the paper's efficacy evidence.
